% from ega4-16-fr.tex
\oldpage[IV-4]{53}

La question étant locale, on peut se borner au cas où $S$ et $X$ sont
affines, et par suite le sous-préschéma diagonal de $X\times_S X$ est fermé.
L'hypothèse que $f$ est localement de type fini entraîne que $\Delta_f$ est
localement de présentation finie (I, 4.3.1), donc le sous-préschéma diagonal
de $X\times_S X$ est défini par un Idéal de type fini $\mathscr{I}$, et
$\Omega_{X/S}^1=\mathscr{I}/\mathscr{I}^2$ est un $\mathcal{O}_X$-Module de
type fini. La proposition résulte alors aussitôt de la comparaison des
conditions dans (16.10.1) et (16.9.4).

\begin{remark}[16.10.6]
\label{IV.16.10.6}
Soit $f:X\to S$ un morphisme tel que le $\mathcal{O}_X$-Module
$\Omega_{X/S}^1$ soit localement libre de rang fini. Il résulte alors de
(0, 20.4.7) que tout $x\in X$ possède un voisinage ouvert $U$ tel qu'il existe
une famille finie $(z_\lambda)_{\lambda\in L}$ de sections de
$\mathcal{O}_X$ au-dessus de $U$ pour lesquelles $(dz_\lambda)_{\lambda\in L}$
forme une \emph{base} du $\Gamma(U,\mathcal{O}_X)$-module
$\Gamma(U,\Omega_{X/S}^1)$.
\end{remark}

\subsection{\texorpdfstring{Opérateurs différentiels sur un $S$-préschéma différentiellement
lisse}{Opérateurs différentiels sur un S-préschéma différentiellement lisse}}
\EGAReaderTarget{ega4.c04.subsec.16-11}{IV.16.11}

\begin{env}[16.11.1]
\label{IV.16.11.1}
Soient $f:X\to S$ un morphisme, $U$ un ouvert de $X$,
$(z_\lambda)_{\lambda\in L}$ une famille de sections de $\mathcal{O}_X$
au-dessus de $U$ telle que les $dz_\lambda$ forment un système de générateurs
de $\Omega_{X/S}^1|U=\Omega_{U/S}^1$. Soit $m$ un entier ou le symbole
$\infty$, et posons, pour tout $\lambda$
\[
  \zeta_\lambda=\delta z_\lambda=d^m z_\lambda-z_\lambda
  \in\Gamma(U,\mathscr{P}_{X/S}^m).
  \tag{16.11.1.1}
  \label{IV.16.11.1.1}
\]

Nous utiliserons d'autre part les notations usuelles de l'analyse; pour tout
$\mathbf{p}=(p_\lambda)\in\mathbb{N}^{(L)}$ (avec $p_\lambda=0$ sauf pour un
nombre fini d'indices), nous poserons
\[
  |\mathbf{p}|=\sum_\lambda p_\lambda,
  \qquad \mathbf{p}!=\prod_\lambda(p_\lambda!)
  \tag{16.11.1.2}
  \label{IV.16.11.1.2}
\]
\[
  \binom{\mathbf{p}}{\mathbf{q}}
  =\mathbf{p}!/(\mathbf{q}!(\mathbf{p}-\mathbf{q})!)
  \quad\text{pour $\mathbf{p},\mathbf{q}$ dans $\mathbb{N}^{(L)}$,
  $\mathbf{q}\leq\mathbf{p}$}
  \tag{16.11.1.3}
  \label{IV.16.11.1.3}
\]
et on convient que $\binom{\mathbf{p}}{\mathbf{q}}=0$ si
$\mathbf{q}\nleq\mathbf{p}$,
\[
  \mathbf{z}^{\mathbf{p}}=\prod_\lambda(z_\lambda)^{p_\lambda},
  \qquad
  \boldsymbol{\zeta}^{\mathbf{p}}=\prod_\lambda(\zeta_\lambda)^{p_\lambda}.
  \tag{16.11.1.4}
  \label{IV.16.11.1.4}
\]

On aura donc, avec ces notations
\[
  d^m(\mathbf{z}^{\mathbf{p}})
  =(d^m(\mathbf{z}))^{\mathbf{p}}
  =(\boldsymbol{\zeta}+\mathbf{z})^{\mathbf{p}}
  =\sum_{\mathbf{q}\leq\mathbf{p}}
    \binom{\mathbf{p}}{\mathbf{q}}
    \mathbf{z}^{\mathbf{p}-\mathbf{q}}\boldsymbol{\zeta}^{\mathbf{q}}
  \tag{16.11.1.5}
  \label{IV.16.11.1.5}
\]
\[
  \boldsymbol{\zeta}^{\mathbf{p}}
  =(d^m\mathbf{z}-\mathbf{z})^{\mathbf{p}}
  =\sum_{\mathbf{q}\leq\mathbf{p}}
    (-1)^{|\mathbf{p}-\mathbf{q}|}
    \binom{\mathbf{p}}{\mathbf{q}}
    \mathbf{z}^{\mathbf{p}-\mathbf{q}}d^m(\mathbf{z}^{\mathbf{q}}).
  \tag{16.11.1.6}
  \label{IV.16.11.1.6}
\]

Comme les $dz_\lambda$ engendrent $\Omega_{X/S}^1$ et sont les images des
$\delta z_\lambda$, et que l'homomorphisme canonique (16.3.1.1) est
surjectif, on en conclut que pour $m$ fini, les $\delta z_\lambda$ engendrent
la $\mathcal{O}_U$-Algèbre $\mathscr{P}_{U/S}^m$ (Bourbaki,
\emph{Alg. comm.}, chap.~III, §~2, n\textsuperscript{o}~8, cor.~2 du th.~1).
Donc les $\boldsymbol{\zeta}^{\mathbf{p}}$ (pour $|\mathbf{p}|\leq m$)
engendrent le $\mathcal{O}_U$-Module $\mathscr{P}_{U/S}^m$. Un opérateur
différentiel $D\in\operatorname{Diff}_{U/S}^m$ est par suite entièrement
déterminé par les valeurs des
$\langle\boldsymbol{\zeta}^{\mathbf{p}},D\rangle$ pour
$|\mathbf{p}|\leq m$, ou, ce qui revient au même par (16.11.1.5) et
(16.11.1.6), par les valeurs
